% !TeX root = ../main.tex
\chapter{Théorie}
\label{chap:theorie}

\setlength{\abovedisplayskip}{5pt}
\setlength{\belowdisplayskip}{5pt}
\setlength{\abovedisplayshortskip}{2pt}
\setlength{\belowdisplayshortskip}{2pt}

Ce chapitre présente les équations utiles au traitement des mesures : le comportement électrique des capteurs (CTN et Pt100), le conditionnement en tension, le montage 4 fils Kelvin, le modèle dynamique du premier ordre et le filtrage par suréchantillonnage.

\section{Thermistance CTN}
\label{sec:theorie_ntc}

La thermistance CTN (\textit{Negative Temperature Coefficient}) est un semi-conducteur dont la résistance diminue fortement avec la température. Pour de petits intervalles autour de l'ambiance, la température s'obtient par régression linéaire locale :
\begin{equation}
    T = T_{\text{ref}} + \frac{R(T_{\text{ref}}) - R_{\text{mes}}}{p_{\text{moy}}}
    \label{eq:theorie_pente_locale}
\end{equation}
où $p_{\text{moy}} = -\frac{\Delta R}{\Delta T}$ [$\unit{\kilo\ohm\per\degreeCelsius}$] est la sensibilité locale issue de la datasheet constructeur.

Pour mesurer la résistance avec la carte d'acquisition, la thermistance est insérée dans un pont diviseur alimenté sous $V_{CC}$ avec une résistance série $R$ :
\begin{equation}
    U_{NTC} = V_{CC} \cdot \frac{R_{NTC}}{R + R_{NTC}} \quad \iff \quad R_{NTC} = \frac{R \cdot U_{NTC}}{V_{CC} - U_{NTC}}
    \label{eq:theorie_pont}
\end{equation}
Le courant de mesure dissipe dans la perle une puissance Joule $P$ causant un auto-échauffement interne $\Delta T_{\text{auto}}$ :
\begin{equation}
    P = \frac{U_{NTC}^2}{R_{NTC}} \quad \text{et} \quad \Delta T_{\text{auto}} = \frac{P}{\delta}
    \label{eq:theorie_puissance_auto}
\end{equation}
où $\delta$ est le facteur de dissipation en air calme [$\unit{\milli\watt\per\degreeCelsius}$]. L'échauffement $\Delta T_{\text{auto}}$ doit rester négligeable devant la tolérance du capteur ($\pm \qty{0.8}{\degreeCelsius}$).

\section{Sonde thermorésistive Pt100}
\label{sec:theorie_pt100}

La sonde Pt100 utilise un élément en platine pur de résistance normalisée $R_0 = \qty{100.00}{\ohm}$ à \qty{0}{\degreeCelsius}. Entre \qty{0}{\degreeCelsius} et \qty{100}{\degreeCelsius}, sa caractéristique est modélisée par la relation linéaire (IEC 60751) :
\begin{equation}
    R(T) = R_0 \cdot (1 + \alpha\,T) \quad \text{avec } \alpha = \qty{3.850e-3}{\degreeCelsius^{-1}}
    \label{eq:theorie_pt100_lineaire}
\end{equation}
La sensibilité thermique nominale et la formule inverse de température valent :
\begin{equation}
    S = R_0 \cdot \alpha = 100.00 \times 3.850 \cdot 10^{-3} = \qty{0.3850}{\ohm\per\degreeCelsius} \quad \implies \quad T = \frac{R(T) - R_0}{S}
    \label{eq:theorie_sensibilite_pt100}
\end{equation}

La résistance d'une Pt100 étant faible (\qty{100}{\ohm}), les conducteurs de liaison ont un impact direct sur la mesure :
\begin{itemize}[leftmargin=1.2cm]
    \item \textbf{Montage 2 fils (2W) :} Le multimètre mesure la résistance totale incluant les deux fils de ligne ($2 R_{\text{ligne}}$), introduisant un biais systématique positif :
    \begin{equation}
        R_{\text{mes, 2W}} = R_{\text{Pt100}} + 2 \cdot R_{\text{ligne}} \quad \implies \quad \Delta T = \frac{2 \cdot R_{\text{ligne}}}{S}
    \end{equation}
    \item \textbf{Montage 4 fils (4W) Kelvin :} Deux conducteurs injectent un courant constant ($I_{\text{inj}} \le \qty{1}{\milli\ampere}$), tandis que deux autres mesurent la tension aux bornes de la sonde. L'impédance du voltmètre étant très grande ($> \qty{10}{\giga\ohm}$), le courant dans les fils de tension est quasi nul ($I_{\text{sense}} \approx 0$). Aucune chute de tension ne s'y produit : l'influence de la résistance des câbles est rigoureusement annulée.
\end{itemize}

\section{Modèle dynamique du premier ordre}
\label{sec:theorie_dynamique}

Lorsqu'une sonde est plongée dans un fluide à température $T_\infty$, le bilan convectif s'écrit selon la loi de Newton :
\begin{equation}
    m \cdot c \cdot \frac{\mathrm{d}T(t)}{\mathrm{d}t} = h \cdot S \cdot \big(T_\infty - T(t)\big) \quad \iff \quad \tau \cdot \frac{\mathrm{d}T(t)}{\mathrm{d}t} + T(t) = T_\infty
    \label{eq:theorie_equa_diff}
\end{equation}
avec la constante de temps thermique $\tau = \frac{m \cdot c}{h \cdot S}$ [s] ($m$ masse, $c$ chaleur massique, $h$ coefficient convectif, $S$ surface).

Pour une température initiale $T(0) = T_0$, la solution indicielle s'écrit :
\begin{equation}
    T(t) = T_\infty - (T_\infty - T_0) \cdot e^{-t/\tau}
    \label{eq:theorie_solution_indicielle}
\end{equation}
Quatre méthodes sont utilisées pour identifier la constante de temps :
\begin{enumerate}[leftmargin=1.2cm]
    \item \textbf{Tangente à l'origine :} La tangente en $t = 0$ intercepte l'asymptote $T = T_\infty$ à $t = \tau$.
    \item \textbf{Seuil à \qty{63.2}{\percent} :} Pour $t = \tau$, le capteur a parcouru $1 - e^{-1} \approx \qty{63.2}{\percent}$ du saut thermique.
    \item \textbf{Fit exponentiel :} Minimisation de la somme des carrés des résidus du modèle $T(t) = A + B \cdot e^{-t/\tau}$.
    \item \textbf{Temps de réponse à \qty{95}{\percent} ($t_{95\%}$) :} Temps pour atteindre \qty{95}{\percent} du saut : $t_{95\%} = -\ln(0.05) \cdot \tau \approx 3 \cdot \tau$.
\end{enumerate}

\section{Réduction du bruit par suréchantillonnage}
\label{sec:theorie_oversampling}

Pour un bruit blanc d'écart-type $\sigma$, la moyenne arithmétique de $N$ échantillons présente un écart-type réduit par $\sqrt{N}$ :
\begin{equation}
    \sigma_{\bar{x}} = \frac{\sigma}{\sqrt{N}}
    \label{eq:theorie_reduction_bruit}
\end{equation}
Pour $N = 100$ échantillons sous LabVIEW, le bruit électronique de conversion est théoriquement divisé par $\sqrt{100} = 10$.
